%%%%%%%%%%%%%%%%%%%%%%%%%%%%%%%%%%%%%%%%%%%%%%%%%%%%%%%%%%%%%%%%%%%%%%
% Lớp 9 | Chương V. Đường tròn | Mở đầu về đường tròn
% Dạng 9H5?13-1: Vị trí của điểm đối với đường tròn
%%%%%%%%%%%%%%%%%%%%%%%%%%%%%%%%%%%%%%%%%%%%%%%%%%%%%%%%%%%%%%%%%%%%%%


% nguồn: dan_tu_AI (nhập 14/11/2026)
\begin{ex}%[9H5N13-1]%[ma:c5dd1421]
	Cho đường tròn $(O;5\text{ cm})$ và điểm $M$ sao cho $OM=4\text{ cm}$. Vị trí của điểm $M$ đối với đường tròn $(O)$ là
	\choice
	{\True $M$ nằm bên trong đường tròn}
	{$M$ nằm trên đường tròn}
	{$M$ nằm bên ngoài đường tròn}
	{$M$ trùng với tâm $O$}
	\loigiai{
		Vì $OM=4<5=R$ nên $M$ nằm bên trong đường tròn $(O)$.
	}
\end{ex}

% nguồn: dan_tu_AI (nhập 14/11/2026)
\begin{ex}%[9H5H13-1]%[ma:1c0af5cb]
	Cho tam giác $ABC$ vuông tại $A$ có $AB=6\text{ cm}$, $AC=8\text{ cm}$. Vị trí của điểm $C$ đối với đường tròn $(B;9\text{ cm})$ là
	\choice
	{$C$ nằm bên trong đường tròn}
	{$C$ nằm trên đường tròn}
	{$C$ trùng với tâm đường tròn}
	{\True $C$ nằm bên ngoài đường tròn}
	\loigiai{
		Theo định lí Pythagore: $BC=\sqrt{6^2+8^2}=10\text{ (cm)}$.\\
		Vì $BC=10>9$ nên $C$ nằm bên ngoài đường tròn $(B;9\text{ cm})$.
	}
\end{ex}

% nguồn: dan_tu_AI (nhập 14/11/2026)
\begin{bt}%[9H5V13-1]%[ma:d590b8ba]
	Cho hình chữ nhật $ABCD$ có $AB=6\text{ cm}$, $BC=8\text{ cm}$. Có bao nhiêu giá trị nguyên của $R$ (đơn vị: cm) để điểm $B$ nằm bên trong và điểm $C$ nằm bên ngoài đường tròn $(A;R)$?
	\shortans{3}
	\loigiai{
		$AC=\sqrt{AB^2+BC^2}=\sqrt{36+64}=10\text{ (cm)}$.\\
		$B$ nằm bên trong $(A;R)\Leftrightarrow AB<R\Leftrightarrow R>6$.\\
		$C$ nằm bên ngoài $(A;R)\Leftrightarrow AC>R\Leftrightarrow R<10$.\\
		Vậy $6<R<10$, có $3$ giá trị nguyên là $7$; $8$; $9$.
	}
\end{bt}

% nguồn: dan_tu_AI (nhập 14/11/2026)
\begin{bt}%[9H5V13-1]%[ma:56cc255b]
	Một trạm phát sóng đặt tại $O$ phủ sóng trong phạm vi hình tròn bán kính $3\text{ km}$. Nhà bạn An ở vị trí $A$ cách $O$ là $2\text{ km}$ về phía Đông; nhà bạn Bình ở vị trí $B$ cách $A$ là $2{,}5\text{ km}$ về phía Bắc. Nhà Bình chưa được phủ sóng. Hỏi nếu đi trên đường thẳng từ $B$ về phía $A$ thì phải đi ít nhất bao nhiêu mét mới vào vùng phủ sóng (làm tròn đến hàng đơn vị)?
	\shortans{264}
	\loigiai{
		Tam giác $OAB$ vuông tại $A$: $OB=\sqrt{2^2+2{,}5^2}=\sqrt{10{,}25}>3$ nên $B$ ở ngoài vùng phủ sóng.\\
		Gọi $M$ là điểm trên đoạn $AB$ với $OM=3\text{ km}$. Khi đó $AM=\sqrt{3^2-2^2}=\sqrt5\approx 2{,}236\text{ (km)}$.\\
		Các điểm $N$ trên đoạn $AB$ có $AN\le AM$ thì $ON\le 3$ (nằm trong vùng phủ sóng).\\
		Quãng đường cần đi: $BM=2{,}5-\sqrt5\approx 0{,}264\text{ (km)}\approx 264\text{ (m)}$.
	}
\end{bt}

% nguồn: dan_tu_AI (nhập 18/11/2026)
\begin{ex}%[9H5N13-1]%[ma:d444a530]
	Cho đường tròn $(O; R)$ và một điểm $M$. Khẳng định nào sau đây đúng?
	\choice
	{\True Nếu $OM < R$ thì điểm $M$ nằm trong đường tròn}
	{Nếu $OM > R$ thì điểm $M$ nằm trên đường tròn}
	{Nếu $OM = R$ thì điểm $M$ nằm ngoài đường tròn}
	{Nếu $OM \le R$ thì điểm $M$ nằm ngoài đường tròn}
	\loigiai{
		Theo định nghĩa vị trí tương đối của điểm và đường tròn, điểm $M$ nằm trong đường tròn khi khoảng cách từ tâm đến điểm đó nhỏ hơn bán kính ($OM < R$).
	}
\end{ex}

% nguồn: dan_tu_AI (nhập 18/11/2026)
\begin{ex}%[9H5N13-1]%[ma:af196dba]
	Cho đường tròn tâm $O$, bán kính $R = 5\text{ cm}$. Điểm $A$ cách $O$ một khoảng $4\text{ cm}$. Vị trí của điểm $A$ đối với đường tròn là
	\choice
	{Nằm trên đường tròn}
	{Nằm ngoài đường tròn}
	{\True Nằm trong đường tròn}
	{Trùng với tâm $O$}
	\loigiai{
		Ta có $OA = 4\text{ cm} < R = 5\text{ cm}$ nên điểm $A$ nằm trong đường tròn.
	}
\end{ex}

% nguồn: dan_tu_AI (nhập 18/11/2026)
\begin{ex}%[9H5N13-1]%[ma:2e7f01fb]
	Trong mặt phẳng tọa độ $Oxy$, cho điểm $M(0; -3)$. Điểm $M$ đối với đường tròn $(O; 3)$ sẽ
	\choice
	{Nằm ngoài đường tròn}
	{Nằm trong đường tròn}
	{\True Nằm trên đường tròn}
	{Trùng với tâm đường tròn}
	\loigiai{
		Khoảng cách từ gốc tọa độ $O(0; 0)$ đến $M(0; -3)$ là $OM = \sqrt{0^2 + (-3)^2} = 3$. 
		Vì $OM = R = 3$ nên điểm $M$ nằm trên đường tròn.
	}
\end{ex}

% nguồn: dan_tu_AI (nhập 18/11/2026)
\begin{ex}%[9H5H13-1]%[ma:d088899b]
	Cho tam giác $ABC$ vuông tại $A$ có $AB = 6\text{ cm}, AC = 8\text{ cm}$. Xét đường tròn tâm $B$, bán kính $R = 9\text{ cm}$. Vị trí của điểm $C$ đối với đường tròn $(B; 9)$ là
	\choice
	{Nằm trên đường tròn}
	{Nằm trong đường tròn}
	{Trùng với tâm đường tròn}
	{\True Nằm ngoài đường tròn}
	\loigiai{
		Áp dụng định lí Pythagore trong tam giác vuông $ABC$:
		$BC = \sqrt{AB^2 + AC^2} = \sqrt{6^2 + 8^2} = 10\text{ cm}$.
		Khoảng cách từ tâm $B$ đến $C$ là $10\text{ cm} > 9\text{ cm} = R$. Vậy $C$ nằm ngoài đường tròn.
	}
\end{ex}

% nguồn: dan_tu_AI (nhập 18/11/2026)
\begin{ex}%[9H5H13-1]%[ma:da14e0fe]
	Trong mặt phẳng tọa độ $Oxy$, cho đường tròn tâm $O(0;0)$ bán kính $R = 5$. Điểm $N(-3; 4)$ có vị trí như thế nào đối với đường tròn?
	\choice
	{Nằm ngoài đường tròn}
	{\True Nằm trên đường tròn}
	{Nằm trong đường tròn}
	{Trùng với gốc $O$}
	\loigiai{
		Khoảng cách $ON = \sqrt{(-3)^2 + 4^2} = \sqrt{9 + 16} = 5$.
		Vì $ON = R = 5$ nên điểm $N$ nằm trên đường tròn.
	}
\end{ex}

% nguồn: dan_tu_AI (nhập 18/11/2026)
\begin{ex}%[9H5H13-1]%[ma:d88a23d9]
	Một trạm phát sóng WiFi có vùng phủ sóng là một hình tròn bán kính $15\text{ m}$. Một điện thoại cách trạm phát sóng $12\text{ m}$. Vị trí của điện thoại này so với vùng phủ sóng là
	\choice
	{\True Nằm trong vùng phủ sóng}
	{Nằm trên ranh giới vùng phủ sóng}
	{Nằm ngoài vùng phủ sóng}
	{Không xác định được}
	\loigiai{
		Khoảng cách từ điện thoại đến trạm là $d = 12\text{ m} < R = 15\text{ m}$. Do đó điện thoại nằm trong vùng phủ sóng (nằm trong đường tròn).
	}
\end{ex}

% nguồn: dan_tu_AI (nhập 18/11/2026)
\begin{ex}%[9H5H13-1]%[ma:e296ac00]
	Hai đường tròn đồng tâm $(O; 3\text{ cm})$ và $(O; 5\text{ cm})$ tạo thành một hình vành khăn. Điểm $K$ nằm trong hình vành khăn (không tính các biên) khi khoảng cách $OK$ thỏa mãn
	\choice
	{$OK < 3\text{ cm}$}
	{$OK > 5\text{ cm}$}
	{\True $3\text{ cm} < OK < 5\text{ cm}$}
	{$OK \le 3\text{ cm}$}
	\loigiai{
		Để $K$ nằm trong hình vành khăn (phần không gian giữa hai đường tròn, không lấy biên) thì điểm $K$ phải nằm ngoài đường tròn nhỏ và nằm trong đường tròn lớn. Tức là $3\text{ cm} < OK < 5\text{ cm}$.
	}
\end{ex}

% nguồn: dan_tu_AI (nhập 18/11/2026)
\begin{ex}%[9H5V13-1]%[ma:6892df80]
	Một chiếc cầu vòm tròn bắc qua một con kênh. Biết khoảng cách giữa hai chân cầu trên mặt nước là $AB = 8\text{ m}$. Bán kính của đường tròn chứa vòm cầu là $R = 5\text{ m}$. Chiều cao của vòm cầu so với mặt nước (phần cao nhất) là
	\choice
	{$1\text{ m}$}
	{\True $2\text{ m}$}
	{$3\text{ m}$}
	{$4\text{ m}$}
	\loigiai{
		Gọi tâm đường tròn là $O$. Kẻ $OH \perp AB$ tại $H$, $H$ là trung điểm $AB$ nên $AH = 4\text{ m}$.
		Áp dụng định lí Pythagore cho tam giác vuông $OAH$: $OH = \sqrt{OA^2 - AH^2} = \sqrt{5^2 - 4^2} = 3\text{ m}$.
		Chiều cao vòm cầu là $h = R - OH = 5 - 3 = 2\text{ m}$.
	}
\end{ex}

% nguồn: dan_tu_AI (nhập 18/11/2026)
\begin{ex}%[9H5V13-1]%[ma:2c7aac9e]
	Trên mặt phẳng tọa độ $Oxy$, cho hai điểm $A(3; 4)$ và $B(6; 0)$. Xét đường tròn $(O; 6)$, khẳng định nào sau đây đúng?
	\choice
	{Cả $A$ và $B$ đều nằm trong đường tròn}
	{$A$ nằm trên, $B$ nằm ngoài đường tròn}
	{\True $A$ nằm trong, $B$ nằm trên đường tròn}
	{Cả $A$ và $B$ đều nằm ngoài đường tròn}
	\loigiai{
		Ta có $OA = \sqrt{3^2 + 4^2} = 5 < 6 \Rightarrow A$ nằm trong đường tròn.
		$OB = \sqrt{6^2 + 0^2} = 6 \Rightarrow B$ nằm trên đường tròn.
	}
\end{ex}

% nguồn: dan_tu_AI (nhập 18/11/2026)
\begin{ex}%[9H5C13-1]%[ma:1f3499f1]
	Cho đường tròn $(O; R)$ và một điểm $A$ cố định nằm ngoài đường tròn. Điểm $M$ di động trên đường tròn $(O; R)$. Độ dài đoạn thẳng $AM$ đạt giá trị nhỏ nhất khi và chỉ khi
	\choice
	{$M$ là giao điểm của đường thẳng $OA$ với đường tròn, nằm khác phía $A$ so với $O$}
	{$M$ nằm trên đường thẳng vuông góc với $OA$ tại $O$}
	{$M$ tạo với $O$ và $A$ thành tam giác đều}
	{\True $M$ là giao điểm của đường thẳng $OA$ với đường tròn, nằm giữa $O$ và $A$}
	\loigiai{
		Theo bất đẳng thức tam giác cho $\triangle OAM$, ta luôn có: $AM + OM \ge OA \Rightarrow AM \ge OA - OM = OA - R$.
		Dấu "=" xảy ra khi $M$ nằm giữa $O$ và $A$. Vậy $AM$ nhỏ nhất khi $M$ là giao điểm của đoạn $OA$ với đường tròn.
	}
\end{ex}

% nguồn: dan_tu_AI (nhập 18/11/2026)
\begin{ex}%[9H5H13-1]%[ma:a746cc30]
	Trong mặt phẳng tọa độ $Oxy$, cho đường tròn tâm $O$ bán kính $R = 10$. Các mệnh đề sau đây đúng hay sai?
	\choiceTF
	{Điểm $A(6; 8)$ nằm ngoài đường tròn tâm $O$}
	{\True Điểm $B(-8; -6)$ nằm trên đường tròn tâm $O$}
	{Điểm $C(5; 5)$ nằm trên đường tròn tâm $O$}
	{\True Điểm $D(10; 1)$ nằm ngoài đường tròn tâm $O$}
	\loigiai{
		\begin{itemchoice}
			\itemch Sai. $OA = \sqrt{6^2 + 8^2} = 10 = R \Rightarrow A$ nằm trên đường tròn.
			\itemch Đúng. $OB = \sqrt{(-8)^2 + (-6)^2} = 10 = R \Rightarrow B$ nằm trên đường tròn.
			\itemch Sai. $OC = \sqrt{5^2 + 5^2} = \sqrt{50} < 10 \Rightarrow C$ nằm trong đường tròn.
			\itemch Đúng. $OD = \sqrt{10^2 + 1^2} = \sqrt{101} > 10 \Rightarrow D$ nằm ngoài đường tròn.
		\end{itemchoice}
	}
\end{ex}

% nguồn: dan_tu_AI (nhập 18/11/2026)
\begin{ex}%[9H5V13-1]%[ma:3a572c6b]
	Một tàu cảnh sát biển phát hiện tàu lạ tại tọa độ $M(x; y)$ cách trạm radar gốc $O$ theo hệ thức $x^2 + y^2 = d^2$. Vùng quét của radar là hình tròn tâm $O$ bán kính $20\text{ hải lý}$. Các mệnh đề sau đúng hay sai?
	\choiceTF
	{\True Nếu $d = 15$, tàu lạ nằm trong vùng quét của radar}
	{Nếu $d = 20$, tàu lạ nằm ngoài vùng quét của radar}
	{\True Nếu $d = 25$, tàu lạ không bị radar phát hiện}
	{Vùng quét của radar có đường kính là $20\text{ hải lý}$}
	\loigiai{
		\begin{itemchoice}
			\itemch Đúng. Khoảng cách $OM = d = 15 < 20$, nằm trong vùng quét.
			\itemch Sai. $d = 20 = R$, tàu nằm trên biên vùng quét (vẫn được phát hiện).
			\itemch Đúng. $d = 25 > 20$, nằm ngoài vùng quét nên không bị phát hiện.
			\itemch Sai. Bán kính là $20$ nên đường kính là $40\text{ hải lý}$.
		\end{itemchoice}
	}
\end{ex}

% nguồn: dan_tu_AI (nhập 18/11/2026)
\begin{bt}%[9H5H13-1]%[ma:90c2f19a]
	Cho đường tròn $(O; 4\text{ cm})$. Trên tia $Ox$, lấy các điểm $A_1, A_2, A_3, A_4, A_5$ sao cho khoảng cách từ $O$ đến $A_n$ bằng $n\text{ cm}$ ($n=1, 2, 3, 4, 5$). Có bao nhiêu điểm nằm trong đường tròn $(O; 4\text{ cm})$?
	\shortans{3}
	\loigiai{
		Điểm nằm trong đường tròn khi khoảng cách đến tâm $< R = 4\text{ cm}$.
		Các điểm thỏa mãn là $A_1, A_2, A_3$ (có khoảng cách là 1, 2, 3 cm). Có 3 điểm. ($A_4$ nằm trên, $A_5$ nằm ngoài).
	}
\end{bt}

% nguồn: dan_tu_AI (nhập 18/11/2026)
\begin{bt}%[9H5V13-1]%[ma:e72d16d8]
	Cho đường tròn tâm $O$, bán kính $R = 10\text{ cm}$ và điểm $A$ nằm ngoài đường tròn sao cho $OA = 26\text{ cm}$. Một điểm $M$ di động trên đường tròn. Tính khoảng cách lớn nhất từ $A$ đến $M$ (theo đơn vị cm).
	\shortans{36}
	\loigiai{
		Theo bất đẳng thức tam giác cho ba điểm $O, A, M$, ta có $AM \le AO + OM = 26 + 10 = 36\text{ cm}$.
		Dấu "=" xảy ra khi $O$ nằm giữa $A$ và $M$. Khoảng cách lớn nhất là $36\text{ cm}$.
	}
\end{bt}

% nguồn: dan_tu_AI (nhập 18/11/2026)
\begin{bt}%[9H5C13-1]%[ma:ce4e7b60]
	Trong mặt phẳng tọa độ $Oxy$, có bao nhiêu điểm có tọa độ nguyên $(x; y)$ nằm đúng trên đường tròn tâm $O$ bán kính $R = 5$?
	\shortans{12}
	\loigiai{
		Ta cần tìm số nghiệm nguyên của phương trình $x^2 + y^2 = 25$.
		Các bộ số nguyên là: $(\pm 5; 0), (0; \pm 5), (\pm 3; \pm 4), (\pm 4; \pm 3)$.
		Có $2 + 2 + 4 + 4 = 12$ điểm.
	}
\end{bt}

% nguồn: dan_tu_AI (nhập 18/11/2026)
\begin{bt}%[9H5H13-1]%[ma:4fd5b3dd]
	Một ngọn hải đăng được đặt tại vị trí $O$ chiếu sáng một vùng biển là hình tròn tâm $O$ bán kính $18\text{ hải lý}$. Hai chiếc tàu $A$ và $B$ lần lượt cách ngọn hải đăng các khoảng cách là $15\text{ hải lý}$ và $20\text{ hải lý}$. 
	a) Tàu nào nằm trong vùng nhận được ánh sáng từ ngọn hải đăng? Tàu nào không nhận được?
	b) Giải thích vì sao ánh sáng hải đăng không có trục đối xứng duy nhất (về mặt hình học mặt sáng)?
	\loigiai{
		a) Vùng nhận được ánh sáng là hình tròn $(O; 18)$.
		Ta có $OA = 15\text{ hải lý} < 18\text{ hải lý} \Rightarrow$ Tàu $A$ nằm trong vùng chiếu sáng.
		$OB = 20\text{ hải lý} > 18\text{ hải lý} \Rightarrow$ Tàu $B$ nằm ngoài vùng chiếu sáng, không nhận được ánh sáng. \\
		b) Vùng sáng được mô phỏng là một hình tròn hoàn hảo. 
		Theo tính chất đối xứng của đường tròn, mọi đường thẳng đi qua tâm đều là trục đối xứng. 
		Do đó, vùng sáng có vô số trục đối xứng chứ không phải có duy nhất một trục đối xứng.
	}
\end{bt}
